\documentclass[11pt]{article}
\usepackage{mathblatt}
\usepackage{adjustbox,varwidth}
% gesetzt von werkzeuge/setzer.py; Lösungen nach bau/bauregeln.md „Lösungen“.
\newcommand{\kennung}{FWE}
\begin{document}
\pfheftstil
\fancyfoot[C]{{\footnotesize\color{mbgrau}\kennung\hfill\thepage}}
\pfheftkopf{Streifen- und Kreisdiagramm}{Klasse 7 \textperiodcentered{} Lösungen}

\begin{pfloesung}{}
\lz{1b)}{\mbox{a) $3$ cm} \mbox{b) $4{,}5$ cm} \mbox{c) $0{,}8$ cm} \mbox{d) $6{,}2$ cm}}{a) $30 : 10 = 3$ cm; b) $45 : 10 = 4{,}5$ cm; c) $8 : 10 = 0{,}8$ cm; d) $62 : 10 = 6{,}2$ cm}{}
\end{pfloesung}
\begin{pfloesung}{}
\lz{2.}{Eistee $30\,\%$; $3$\,/\,$2{,}5$\,/\,$1{,}5$\,/\,$3$ cm}{Eistee $= 100\,\% - 30\,\% - 25\,\% - 15\,\% = 30\,\%$; Wasser $= 3$ cm, Apfelsaft $= 2{,}5$ cm, Kakao $= 1{,}5$ cm, Eistee $= 3$ cm\par\smallskip \adjustbox{max width=\linewidth,max height=4cm}{\begin{varwidth}{50cm}\streifenvoll{Wasser/30,Apfelsaft/25,Kakao/15,Eistee/30}\end{varwidth}}}{}
\end{pfloesung}
\begin{pfloesung}{}
\lz{3.}{7a ($25\,\%$ gegen $15\,\%$)}{7a: zu Fuß $= 100\,\% - 45\,\% - 30\,\% = 25\,\%$ (Längen $4{,}5$; $3$; $2{,}5$ cm); 7b: zu Fuß $= 100\,\% - 35\,\% - 50\,\% = 15\,\%$ (Längen $3{,}5$; $5$; $1{,}5$ cm). In der 7a: Ihr Abschnitt „zu Fuß“ ist länger ($2{,}5$ cm gegen $1{,}5$ cm).\par\smallskip \adjustbox{max width=\linewidth,max height=4cm}{\begin{varwidth}{50cm}\streifenvoll{Rad/45,Bus/30,Fuß/25}\par\vspace{2mm}\streifenvoll{Rad/35,Bus/50,Fuß/15}\end{varwidth}}}{}
\end{pfloesung}
\begin{pfloesung}{}
\lz{4b)}{\mbox{a) $90^\circ$} \mbox{b) $180^\circ$} \mbox{c) $36^\circ$} \mbox{d) $270^\circ$}}{a) $0{,}25 \cdot 360^\circ = 90^\circ$; b) $0{,}5 \cdot 360^\circ = 180^\circ$; c) $0{,}1 \cdot 360^\circ = 36^\circ$; d) $0{,}75 \cdot 360^\circ = 270^\circ$}{}
\end{pfloesung}
\begin{pfloesung}{}
\lz{5.}{\mbox{a) $126^\circ$} \mbox{b) $43{,}2^\circ$} \mbox{c) $225^\circ$} \mbox{d) $16{,}2^\circ$}}{a) $0{,}35 \cdot 360^\circ = 126^\circ$; b) $0{,}12 \cdot 360^\circ = 43{,}2^\circ$; c) $0{,}625 \cdot 360^\circ = 225^\circ$; d) $0{,}045 \cdot 360^\circ = 16{,}2^\circ$}{}
\end{pfloesung}
\begin{pfloesung}{}
\lz{6.}{\mbox{a) $135^\circ$} \mbox{b) $12{,}6^\circ$}}{a) $9 : 24 = 0{,}375$; $0{,}375 \cdot 360^\circ = 135^\circ$; b) $14 : 400 = 0{,}035$; $0{,}035 \cdot 360^\circ = 12{,}6^\circ$}{}
\end{pfloesung}
\begin{pfloesung}{}
\lz{7.}{Nein, nur $\approx 9{,}3^\circ$}{Nein: Winkel $= 14 : 540 \cdot 360^\circ \approx 9{,}3^\circ$, weniger als $10^\circ$.}{}
\end{pfloesung}
\begin{pfloesung}{}
\lz{8b)}{Kreis mit $180^\circ$, $120^\circ$, $60^\circ$}{Kontrolle: $180^\circ + 120^\circ + 60^\circ = 360^\circ$; die Sektoren nacheinander ab dem Radius abtragen.\par\smallskip \adjustbox{max width=\linewidth,max height=4cm}{\begin{varwidth}{50cm}\kreisdiagramm[1.4]{/180,/120,/60}\end{varwidth}}}{}
\end{pfloesung}
\begin{pfloesung}{}
\lz{9.}{Wasser $151^\circ$, Saft $119^\circ$, Tee $90^\circ$}{Tee $= 100\,\% - 42\,\% - 33\,\% = 25\,\%$; Wasser $0{,}42 \cdot 360^\circ = 151{,}2^\circ \approx 151^\circ$; Saft $0{,}33 \cdot 360^\circ = 118{,}8^\circ \approx 119^\circ$; Tee $0{,}25 \cdot 360^\circ = 90^\circ$; Kontrolle $151^\circ + 119^\circ + 90^\circ = 360^\circ$\par\smallskip \adjustbox{max width=\linewidth,max height=4cm}{\begin{varwidth}{50cm}\kreisdiagramm[1.4]{Wasser/151,Saft/119,Tee/90}\end{varwidth}}}{}
\end{pfloesung}
\begin{pfloesung}{}
\lz{10.}{$150^\circ$, $90^\circ$, $60^\circ$, $60^\circ$}{Licht und Ton $= 36 - 15 - 9 - 6 = 6$; Spielen $15 : 36 \cdot 360^\circ = 150^\circ$; Bühne $9 : 36 \cdot 360^\circ = 90^\circ$; Kostüme $6 : 36 \cdot 360^\circ = 60^\circ$; Licht und Ton $6 : 36 \cdot 360^\circ = 60^\circ$; Kontrolle $150^\circ + 90^\circ + 60^\circ + 60^\circ = 360^\circ$\par\smallskip \adjustbox{max width=\linewidth,max height=4cm}{\begin{varwidth}{50cm}\kreisdiagramm[1.4]{Spielen/150,Bühne/90,Kostüme/60,Licht/60}\end{varwidth}}}{}
\end{pfloesung}
\begin{pfloesung}{}
\lz{11b)}{Halbkreis: Freizeit; mittlerer: Essen; kleinster: Sparen}{Der Halbkreis ist Freizeit, der mittlere Sektor Essen, der kleinste Sparen.}{}
\end{pfloesung}
\begin{pfloesung}{}
\lz{12.}{A Viertel, B Hälfte, C und D weniger; C und D $25\,\%$}{A ein Viertel, B die Hälfte, C und D weniger als ein Viertel; C und D $= 100\,\% - 50\,\% - 25\,\% = 25\,\%$}{}
\end{pfloesung}
\begin{pfloesung}{}
\lz{13.}{a) C Toilette, A Wäsche, D Geschirr, B Kochen und Trinken; b) $22{,}5^\circ$}{a) C Toilette, A Wäsche, D Geschirr, B Kochen und Trinken; b) Winkel $= 30 : 480 \cdot 360^\circ = 22{,}5^\circ$}{}
\end{pfloesung}
\begin{pfloesung}{}
\lz{14.}{Freunde $20\,\%$; $2$\,/\,$1{,}5$\,/\,$4{,}5$\,/\,$2$ cm}{Freunde $= 100\,\% - 20\,\% - 15\,\% - 45\,\% = 20\,\%$; Sport $= 2$ cm, Musik $= 1{,}5$ cm, Medien $= 4{,}5$ cm, Freunde $= 2$ cm\par\smallskip \adjustbox{max width=\linewidth,max height=4cm}{\begin{varwidth}{50cm}\streifenvoll{Sport/20,Musik/15,Medien/45,Freunde/20}\end{varwidth}}}{}
\end{pfloesung}
\begin{pfloesung}{}
\lz{15.}{$64{,}8^\circ$}{Winkel $= 0{,}18 \cdot 360^\circ = 64{,}8^\circ$ (die $650$ braucht man nicht)}{}
\end{pfloesung}
\begin{pfloesung}{}
\lz{16.}{Meer $167^\circ$, Berge $99^\circ$, Stadt $94^\circ$}{Meer $0{,}465 \cdot 360^\circ = 167{,}4^\circ \approx 167^\circ$; Berge $0{,}275 \cdot 360^\circ = 99^\circ$; Stadt $0{,}26 \cdot 360^\circ = 93{,}6^\circ \approx 94^\circ$; Kontrolle $167^\circ + 99^\circ + 94^\circ = 360^\circ$\par\smallskip \adjustbox{max width=\linewidth,max height=4cm}{\begin{varwidth}{50cm}\kreisdiagramm[1.4]{Meer/167,Berge/99,Stadt/94}\end{varwidth}}}{}
\end{pfloesung}
\begin{pfloesung}{}
\lz{17.}{a) B Getränke, D Spiele, A Tombola, C Bastelstand; b) Ja, $825 > 625$}{a) B Getränke, D Spiele, A Tombola, C Bastelstand; b) Ja: $450 + 375 = 825$ €, die Hälfte ist $1\,250 : 2 = 625$ €.}{}
\end{pfloesung}
\end{document}
